
\begin{result}{E.2.19}{Ring-class restrictions and a fixed-field local-sign example}{Proved}
Every \(\sigma\in\Sigma\) of E.2.18 fixes every imaginary-quadratic
ring class field pointwise. There nevertheless exists an actual
\(\sigma\in\Gamma_{\Q^{\mathrm{cm}}}\), nontrivial on
\(\Q(i)^{\mathrm{ab}}\), with an integral tensor comparison on the
following precise diagram: all CM elliptic curves over \(\Qb\)
with geometric endomorphism algebra \(\Q(i)\), all actual abelian
variety lattices isogenous to products of those curves, their
tensor constructions and specified summands, algebraic
correspondences, and all Artin objects with every finite-model
module map. The same construction respects labelled
absolute-Hodge tensors in the abelian version.
Its common scalar is
\[
 \lambda_5=-1,\qquad \lambda_\ell=1\quad(\ell\ne5),
 \tag{E.2.19a}
\]
which is not one global sign. No membership of this \(\sigma\)
in the all-abelian group \(\Sigma\) is asserted.
\end{result}
\begin{proof}
Let \(R=K(j(E))\) be any imaginary-quadratic ring class field.
Apply the finite separating diagram of E.2.2 over \(R\), regard
its objects as underlying-\(\Q\) schemes, and include its
Artin algebra \(A_R\) and module maps. E.2.3 sends the identity
embedding block to the \(\sigma|_R\) block. E.2.2 forces this
restriction to be the identity or actual conjugation. Since
\(\sigma\) fixes \(K\subset\Q^{\mathrm{cm}}\), the latter is
excluded. This proves the first assertion by the explicit
Artin-block maps.

For every imaginary-quadratic \(K\), one has
\[
 K^{\mathrm{ab}}\cap\Q^{\mathrm{cm}}=\Q^{\mathrm{ab}}.
 \tag{E.2.19b}
\]
The inclusion from right to left follows because \(K/\Q\)
is abelian and cyclotomic fields are CM. For the reverse
inclusion, the intersection is Galois over \(\Q\).
The group \(\Gal(K^{\mathrm{ab}}/\Q)\) is generated by its
abelian index-two subgroup \(\Gal(K^{\mathrm{ab}}/K)\) and
actual complex conjugation \(c\), which has order two.
On the intersection the image of \(c\) is central, because
it is central on \(\Q^{\mathrm{cm}}\). The resulting quotient
group is therefore abelian, proving (E.2.19b).

Now take \(K=\Q(i)\), and let \(s\) be the scalar finite
\(K\)-idele (E.2.19a), with infinite component \(1\).
Its norm to \(\Q\) is exactly \(1\), since each local
scalar norm is \(\lambda_\ell^2\). Put
\[
 g_K=\operatorname{art}_K(s)\in\Gal(K^{\mathrm{ab}}/K).
\]
The Artin norm/restriction diagram makes \(g_K\) trivial
on \(\Q^{\mathrm{ab}}\)
\cite[Chapter V, Proposition 5.2 and its following diagram,
pp.~178--179]{GenMilneCFT2020}.
By (E.2.19b), it is compatible with the identity on
\(\Q^{\mathrm{cm}}\). These two maps define an automorphism
of their compositum and extend to \(\Qb\).
One particular extension is specified by the fixed finite
Galois exhaustion of E.2.17: recursively take the first
restriction whose cylinder meets the nonempty compact set
of these compatible extensions. This uses no computability
claim.

The restriction \(g_K\) is nontrivial. Modulo the ray
modulus \(15\), the idele \(s\) represents the rational
residue \(4\): it is \(1\) modulo \(3\) and \(-1\) modulo
\(5\). The ray-class exact sequence has kernel the image
of \(\Z[i]^\times=\{1,-1,i,-i\}\). Only \(1\) among these
units is \(1\) modulo \(3\), and it is not \(-1\) modulo
\(5\). Thus this ray class is nontrivial
\cite[Chapter V, Theorem 1.7 and proof,
pp.~150--151]{GenMilneCFT2020}.
In particular \(\sigma\notin\Gamma_{F_{\mathrm{cm}}}\).
For every order \(\mathcal O_f=\Z+f\Z[i]\), however,
\(s\in\widehat{\mathcal O_f}^{\,\times}\); hence its ring
class is trivial. This is also the rational-congruence
subgroup in the ideal quotient of
\cite[Theorem 9.18, pp.~190--192]{GenCox1989}.

Let \(A/\Qb\) be any elliptic curve of this fixed CM field.
The two CM types have reflex field \(K\); their reflex
norms are the identity and conjugation, both of which
send \(s\) to the same \(\lambda\).
The exact main theorem
\cite[Theorem 9.10, p.~78]{GenMilneCM2020} therefore gives
a \(K\)-linear quasi-isogeny
\[
 \alpha_A:A\longrightarrow\sigma A,\qquad
 \alpha_A(\lambda x)=\sigma x,
 \qquad
 \widehat\alpha_A=\lambda\rho_A(\sigma).
 \tag{E.2.19c}
\]
It is an actual isomorphism. Indeed, write
\(\alpha_A=f/m\), with \(f\) an honest homomorphism.
Integrality on every Tate lattice makes \(f\) vanish
on \(A[m]\); it factors through \([m]\), proving that
\(\alpha_A\) is honest. Apply the same argument to
its inverse. The formula is independent of the labelling
of \(K\), because both types give the same scalar and
completed homology is faithful on homomorphisms.

For every homomorphism \(f:A\to B\), the equality
\(\alpha_Bf=(\sigma f)\alpha_A\) follows after completion
from Galois naturality and the common scalar, and hence
holds exactly. Products and isogenies therefore extend
the construction to every actual lattice isogenous to
a product of these elliptic curves. More explicitly,
transport the product quasi-isogeny through any chosen
isogeny; its completion is again \(\lambda\rho\).
The preceding integrality argument makes it an actual
isomorphism and proves independence of that choice.

On cohomology use inverse duals and exterior powers.
The completed factor in degree \(r\) is \(\lambda^r\).
Correspondences change degree by an even integer, so
their two factors coincide. Their étale cycle-class
equations descend faithfully to the required Betti
equations. This also treats tensors, duals, the stated
projector lattices and labelled absolute-Hodge tensors.
The polarization multiplier is \(1\), since both the
cyclotomic character and the idele norm are \(1\)
\cite[Remark 9.11(c), p.~78]{GenMilneCM2020}.
The Tate component is the identity.

Finally an object defined over an arbitrary finite
field \(L\) has underlying-\(\Q\) Betti realization
\(\bigoplus_u H_B(uA)\). Every \(uA\) still has the
same imaginary-quadratic CM field, though its type
may be conjugated. Both types have already been
treated. Use \(\alpha_{uA}\) from the \(u\)-block to
the \(\sigma u\)-block and put \(e_u\mapsto e_{\sigma u}\)
on every Artin algebra. Naturality proves the module
equations, preserves \(\sum e_u\otimes e_u\), and
commutes with every change of finite model.
The exact descent of E.2.3 gives the asserted global
comparison on this specified fixed-field diagram.
\end{proof}

\begin{result}{E.2.20}{Two CM types eliminate the local-sign ambiguity}{Proved}
The homomorphism \(\delta\) in E.2.18 is trivial. More precisely,
an actual comparison on the full geometric CM/Artin diagram
with Artin restriction in \(\Gamma_{\Q^{\mathrm{cm}}}\)
and common scalar \(\lambda\in\prod_\ell\{1,-1\}\)
forces \(\lambda\) to be one global \(1\) or \(-1\).
Consequently E.2.18 becomes the exact criterion
\[
 \begin{split}
 \Sigma
  &=\{\sigma\in\Gamma_{\Q^{\mathrm{cm}}}:
       \rho_A(\sigma)(\Lambda_A)=\Lambda_{\sigma A}
       \text{ for every abelian variety }A/\Qb\},\\
 I_{\mathrm{ab}}&=\Sigma\sqcup c\Sigma.
 \end{split}
 \tag{E.2.20a}
\]
The sign in E.2.19 is already excluded by two CM types
of the explicit degree-eight field
\[
 L_{15}=\Q(\zeta_{15}+\zeta_{15}^{-1},i).
 \tag{E.2.20b}
\]
\end{result}
\begin{proof}
Fix \(m\ge3\), and write
\[
 F=\Q(\zeta_m+\zeta_m^{-1}),\qquad L=F(i).
\]
Then \(L/\Q\) is Galois and CM, with \(L^{\langle c\rangle}=F\).
Choose a CM type \(\Phi\) containing the identity embedding,
and let \(\Phi'\) replace that embedding by \(c\), leaving
the other conjugate pairs unchanged. The CM construction
gives abelian varieties \(A_\Phi,A_{\Phi'}\) with labelled
\(L\)-action and rational homology of rank one over \(L\).
This includes imprimitive types: the polarized torus
\(\C^\Phi/\Phi(\mathcal O_L)\) is supplied by
\cite[Example 2.11]{GenMilneCM2020}, and
\cite[Proposition 7.9 and Corollary 7.10,
pp.~54--55]{GenMilneCM2020} gives a model and its morphisms
over \(\Qb\). No equality of its full endomorphism algebra
with \(L\) is required. The labelled action and every
map are defined over a finite field, so their actual
Artin-selected blocks are available by E.2.3.

Choose one finite \(L\)-idele \(s\) whose Artin image is
\(\sigma|_{L^{\mathrm{ab}}}\). The global Artin map is
surjective, and the infinite ideles of this totally
imaginary field are connected and have trivial Artin
image, so finite ideles suffice.
For a CM type \(\Psi\) of \(L\), its reflex field
\(E_\Psi\) is contained in \(L\). Norm compatibility
identifies the Artin image of
\(\operatorname{Nm}_{L/E_\Psi}(s)\) with
\(\sigma|_{E_\Psi^{\mathrm{ab}}}\); here
\(E_\Psi^{\mathrm{ab}}L/L\) is abelian and contained
in \(L^{\mathrm{ab}}\), which is the precise
restriction map being used.

The exact norm formulas are
\[
 N_{L,\Psi}
  =N_\Psi\circ\operatorname{Nm}_{L/E_\Psi},
 \qquad
 N_{L,\Psi}(s)=\prod_{\psi\in\Psi}\psi^{-1}(s).
 \tag{E.2.20c}
\]
The first is \cite[Proposition 1.23, p.~16]{GenMilneCM2020};
the second is \cite[Proposition 1.26, p.~17]{GenMilneCM2020}
with \(k=E=L\). Remark 1.25 makes these morphisms of
tori and hence identities on finite ideles.
In particular the norm to the reflex field is already
included; there is no extra stabilizer multiplicity.

The main theorem gives an \(L\)-linear quasi-isogeny
\(\alpha_\Psi:A_\Psi\to\sigma A_\Psi\) with
\[
 \widehat\alpha_\Psi
     =\rho_{A_\Psi}(\sigma)N_{L,\Psi}(s)^{-1}.
\]
The Betti comparison \(T_{A_\Psi}\) is also \(L\)-linear
relative to the two labelled actions, because its
Artin-block naturality is
\[
 T_{A_\Psi}\,\iota(a)=(\sigma\iota(a))\,T_{A_\Psi}
 \quad(a\in L).
\]
Thus \(\alpha_\Psi^{-1}T_{A_\Psi}\) is an \(L\)-linear
automorphism of a rank-one \(L\)-vector space, hence
multiplication by an element \(b_\Psi\in L^\times\).
After completion this says
\[
 \lambda N_{L,\Psi}(s)=b_\Psi\in L^\times.
 \tag{E.2.20d}
\]
This centralizer calculation is valid for imprimitive
types and uses no Hodge preservation of the original
integral comparison.

Take the quotient of (E.2.20d) for \(\Phi,\Phi'\).
Only the identity/conjugate pair changes, so
\[
 r=s/c(s)\in L^\times,\qquad r\,c(r)=1.
\]
An explicit Hilbert-90 solution is
\(b=1+r\) when \(r\ne-1\), and \(b=i\) when \(r=-1\);
in both cases \(b/c(b)=r\).
Consequently \(t=s/b\) is fixed by \(c\), and therefore
\(t\in\mathbb A_{F,f}^{\times}\).
Each CM type contains one extension of every embedding
of \(F\), so (E.2.20c) gives
\[
 N_{L,\Psi}(t)=\operatorname{Nm}_{F/\Q}(t).
\]
Dividing (E.2.20d) by the principal element
\(N_{L,\Psi}(b)\) shows that
\(\lambda\operatorname{Nm}_{F/\Q}(t)\in L^\times\).
It is also a scalar rational finite idele. A diagonal
element of \(L\) with this property is fixed by every
element of \(\Gal(L/\Q)\), hence lies in \(\Q^\times\).
For some \(a\in\Q^\times\), we obtain
\[
 \operatorname{Nm}_{F/\Q}(t)=a\lambda.
 \tag{E.2.20e}
\]

The Artin image of the left side restricts trivially
to \(F\), by the norm/restriction diagram.
The finite principal idele \(a\), with infinite
component \(1\), acts on \(\Q^{\mathrm{ab}}\) as
the identity or actual complex conjugation:
insert its missing infinite component and use
global reciprocity. Both actions are trivial on
the totally real field \(F\).
Equation (E.2.20e) therefore makes
\(\operatorname{art}_{\Q}(\lambda)\) trivial on
\(\Q(\zeta_m+\zeta_m^{-1})\). Equivalently,
\[
 \lambda\bmod m\in\{1,-1\}.
\]
This holds for every \(m\). Their real cyclotomic
fields exhaust \(\Q^{\mathrm{ab},+}\), and
\(\Gal(\Q^{\mathrm{ab}}/\Q^{\mathrm{ab},+})
=\{1,c\}\). Cyclotomic reciprocity identifies these
two elements with the two global units \(1,-1\).
Thus \(\lambda\) is global, as asserted.

For E.2.19 the residue modulo \(15\) is \(4\), which
is neither \(1\) nor \(-1\). The two types just
constructed for \(F_{15}(i)\) give an actual finite
obstruction. These are dimension-four abelian
varieties; include a common finite field of
definition, a finite basis of labelled integral
endomorphisms, and the Artin module maps of E.2.3.
The assertion is about these actual objects and
maps, and does not depend on treating an arbitrary
idele as a Galois element fixing \(\Q^{\mathrm{cm}}\).

Finally a global sign preserves every actual lattice.
Substitution into the proved necessity and sufficiency
of E.2.18 gives (E.2.20a), with exactly the two lifts
of E.2.15. Whether \(\Sigma\) contains a nonidentity
element, whether it connects the quartic place-pair,
and whether an abelian comparison extends to the full
pure diagram remain the precise unresolved questions.
\end{proof}
